\documentclass[11pt]{article}
% Avoid an output-directory-dependent trailer ID in deterministic snapshots.
\ifdefined\pdftrailerid\pdftrailerid{}\fi
\usepackage[letterpaper,margin=1in]{geometry}
\usepackage[T1]{fontenc}
\usepackage[utf8]{inputenc}
\usepackage{lmodern}
\usepackage{amsmath,amssymb,amsthm,mathtools}
\usepackage{microtype}
\usepackage{enumitem}
\usepackage{xurl}
\usepackage[hidelinks,unicode,pdfencoding=auto]{hyperref}
\usepackage{bookmark}
\newtheorem{theorem}{Theorem}[section]
\newtheorem{lemma}[theorem]{Lemma}
\newtheorem{proposition}[theorem]{Proposition}
\newtheorem{corollary}[theorem]{Corollary}
\theoremstyle{remark}
\newtheorem{remark}[theorem]{Remark}
\newcommand{\tightlist}{\setlength{\itemsep}{0pt}\setlength{\parskip}{0pt}}
\setlist{itemsep=2pt,topsep=4pt}
\setlength{\emergencystretch}{2em}
\allowdisplaybreaks[1]
\hypersetup{pdftitle={Cycle Universality for a Single Thorp Shuffle},pdfauthor={Christopher D. Long},pdfsubject={Mesoscopic cycle moments, Poisson limits, and a truncated functional Erdos--Turan law}}
\title{Cycle Universality for a Single Thorp Shuffle}
\author{Christopher D. Long}
\date{9 October 2026}
\begin{document}
\maketitle
\begin{abstract}
We study the cycles of one physical Thorp shuffle on $N=2^d$ states,
equivalently a uniformly random binary noncoalescing feedback shift
register. For every fixed number of prescribed cycles, we prove uniform
joint factorial-moment asymptotics when their lengths are at least $4d$
and their total length is at most $N^\beta$, for fixed $\beta<2/3$.
No restriction is imposed on the ratios of the cycle lengths. These
estimates yield a joint functional Gaussian limit for cycle counts and
logarithmic truncated permutation order on every fixed exponent interval
$[0,\alpha]$, $\alpha<2/3$. They also give Poisson approximation on sets
of bounded harmonic mass, a scale-invariant Poisson cycle process, and
Dickman and uniform limits for truncated cycle mass and the largest cycle
below a cutoff. A separate argument gives a quantitative Poisson
approximation for the entire initial cycle-count vector below the
birthday scale. The proof combines weighted cyclic words, empirical-type
entropy, marked switching, longest-last forced closure, and a
middle-first bridge. The feedback function is sampled once and reused.
The full permutation-order central limit theorem is not proved.
\end{abstract}

\section{Introduction and main
results}\label{introduction-and-main-results}

A random permutation can have universal cycle statistics long before its full
law resembles the uniform measure. We study this distinction for one physical
Thorp shuffle on $N=2^d$ positions. In binary coordinates its permutation is
\[
 T_f(x_1,\ldots,x_d)
 =(x_2,\ldots,x_d,x_1\oplus f(x_2,\ldots,x_d)),
 \qquad f:\mathbb F_2^{d-1}\longrightarrow\mathbb F_2,
\]
where all $N/2$ values of $f$ are independent fair bits. This is the binary
noncoalescing feedback-shift-register model. The feedback function is sampled
\emph{once}; following a cycle reuses its values, rather than supplying an
independent table at every iterate.

Thorp introduced the shuffle in his study of nonrandom shuffling and Faro
\cite{Thorp}. For this feedback model, Arratia, Canfield, and Hales
\cite{ACH} proved the Poisson--Dirichlet limit for the ranked macroscopic
cycle lengths. Stark \cite{Stark} proved exact independent-binomial laws
for sufficiently short cycles. Neither a macroscopic weak limit nor
short-cycle independence determines the joint statistics of cycles at
intermediate lengths. The present paper establishes uniform factorial
moments across that intermediate regime, with no comparability assumption
on the prescribed lengths.

Write $C_r=C_r(T_f)$ for the number of $r$-cycles, $u(x)$ for the last
$d-1$ bits of a state, and $\bar x$ for the state obtained by flipping its
first bit. We use $\mathcal L(Y)$ for a law and
$d_{\rm TV}(\mu,\nu)=\sup_A|\mu(A)-\nu(A)|$. Entropies and $\log_2$
are in bits; other logarithms are natural. The notation $(z)_a$ denotes a
falling factorial, and the least common multiple of an empty set is one.

Our main estimates lead to a truncated functional Erd\H{o}s--Tur\'{a}n
law, complementing the classical permutation-order theorem \cite{ET} and
one-shuffle universality results for riffle and shelf models \cite{Long}.
They also give Poisson approximation on sets of bounded harmonic mass,
a scale-invariant Poisson cycle process, and a Dickman law for the mass in
truncated cycles. A separate quantitative estimate treats the full initial
cycle-count vector below the birthday scale. The proof uses weighted
cyclic words, empirical-type entropy, marked feedback switches, and forced
closures. It uses neither an independent-round mixing theorem nor the
macroscopic Poisson--Dirichlet limit.

\subsection{Joint cycle moments}



For a list of positive lengths \(\mathbf r=(r_1,\ldots,r_k)\), put

\[
 L=\sum_{i=1}^k r_i,\qquad
 a_m=\#\{i:r_i=m\},\qquad
 Q_d(\mathbf r)=\left(\prod_{i=1}^k r_i\right)
 \mathbb E\prod_{m\ge1}(C_m)_{a_m},\qquad Q_d(\varnothing)=1.
\]

Fix an integer \(k\ge1\) and a real \(0<\beta<2/3\). Set

\[
 p_\beta=\frac{2}{2-3\beta},\qquad
 P_{k,\beta}(d)=d^{4k+1}+d(d\log_2d)^{kp_\beta}.
\]

\begin{theorem}[Uniform cycle factorial moments]\label{thm:moments}
Fix $k\ge1$ and $0<\beta<2/3$. For all sufficiently large $d$, depending on
$k$ and $\beta$, every list of positive lengths with $L\le N^\beta$ satisfies
\[
 Q_d(\mathbf r)\le1+C_{k,\beta}P_{k,\beta}(d)\frac LN.
 \tag{1.1}\label{eq:1.1}
\]
If additionally $r_i\ge4d$ for every $i$, then
\[
 1-C_k\frac{dL}{N}\le Q_d(\mathbf r)
 \le1+C_{k,\beta}P_{k,\beta}(d)\frac LN.
 \tag{1.2}\label{eq:1.2}
\]
\end{theorem}

Constants are independent of all ratios \(r_i/r_j\). In particular,

\[
 rs\,\mathbb E[C_r(C_s-\mathbf1_{\{r=s\}})]
 =1+O_\beta\!\left(P_{2,\beta}(d)\frac{r+s}{N}\right)
\]

uniformly for \(r,s\ge4d\) and \(r+s\le N^\beta\). Equal lengths are
covered by the falling factorial, so the cycles are required to be
distinct.

The displayed growth factor is deliberately not optimized. The
conclusion is uniform with exponentially small error in \(d\) on every
fixed exponent range.

\subsection{Functional consequence}\label{functional-consequence}

\begin{theorem}[Truncated functional Erd\H{o}s--Tur\'{a}n law]\label{thm:functional}
Fix $0<\alpha<2/3$, put $\ell=\log N$, and define, for $0\le u\le\alpha$,
\[
 K_d(u)=\sum_{r\le\lfloor N^u\rfloor}C_r,
 \qquad
 O_d(u)=\log\operatorname{lcm}\{r\le\lfloor N^u\rfloor:C_r>0\}.
\]
Then, in $D([0,\alpha],\mathbb R^2)$ with its Skorokhod $J_1$ topology,
\[
 \left(
 \frac{K_d(u)-u\ell}{\sqrt\ell},
 \frac{O_d(u)-\tfrac12u^2\ell^2}{\sqrt{\ell^3/3}}
 \right)_{0\le u\le\alpha}
 \ \Longrightarrow\
 \left(W(u),\sqrt3\int_0^u t\,dW(t)\right)_{0\le u\le\alpha}.
 \tag{1.3}\label{eq:1.3}
\]
Here $W$ is one standard Brownian motion, shared by both coordinates. In
particular, for $b=\lfloor N^\alpha\rfloor$,
\[
 \frac{\log\operatorname{lcm}\{r\le b:C_r>0\}
       -\tfrac12(\log b)^2}{\sqrt{(\log b)^3/3}}
 \ \Longrightarrow\ \mathcal N(0,1).
 \tag{1.4}\label{eq:1.4}
\]
\end{theorem}



The proof of the functional consequence, including tightness, is given
in Sections 8--9. The parameters \(k,\beta,\alpha\) are fixed before
\(d\to\infty\); no uniform assertion as \(k\to\infty\) or
\(\beta\uparrow2/3\) is made.

\section{Exact cycle measures and weighted
words}\label{exact-cycle-measures-and-weighted-words}

Choose \(k\) independent uniform roots, with replacement, independently
of \(f\). Let \(E_{\mathbf r}\) be the event that their cycles are
distinct and have the prescribed lengths. Then the finite measure

\[
 \mu_{\mathbf r}=N^k\mathbb P_0(\,\cdot\,,E_{\mathbf r})
\]

has total mass \(Q_d(\mathbf r)\). An unrooted ordered cycle family has
the same weight at every choice of roots on its cycles.

Equivalently, prescribe \(k\) distinguished cyclic binary words with
these lengths. Let \(A_d\) require all their length-\(d\) windows to be
distinct. A context is a length-\((d-1)\) window immediately following
the initial bit of a state. On \(A_d\), a context has at most two
occurrences. If a context occurs twice, the two preceding bits differ
and the two following bits differ, so the feedback requirements are
consistent. If \(m=L-v\), with \(v\) distinct contexts, the family
occurs with probability \(2^{-v}\). Therefore

\[
 Q_d^{(m)}(\mathbf r)
 =2^{m-L}\#\{\text{valid rooted word families with }m\text{ repeated contexts}\}.
 \tag{2.1}\label{eq:2.1}
\]

This identity counts all repeats, internal or between cycles. No
context-simplicity condition is imposed.

\section{A global entropy tail for a fixed number of
words}\label{a-global-entropy-tail-for-a-fixed-number-of-words}

\subsection{Type counting, including short component
words}\label{type-counting-including-short-component-words}

For a cyclic binary word of length \(n\), let \(H_h\) be the entropy of
its empirical \(h\)-block distribution and \(\delta_h=H_h-H_{h-1}\). A
fixed \(h\)-type contains at most

\[
 n\,2^{n\delta_h}
 \tag{3.1}\label{eq:3.1}
\]

rooted words. For \(h\le n\), use the empirical order-\((h-1)\) Markov
chain: every word of the type has probability at least
\(n^{-1}2^{-n\delta_h}\), since its initial block has mass at least
\(1/n\) and omitting wraparound transition factors only increases the
product.

The bound is also valid for \(h>n\), a point needed for widely separated
lengths. An \(h\)-block then contains the entire cyclic word. A fixed
\(h\)-type determines its necklace, so at most \(n\) rooted words have
the type. Also \(\delta_h=0\) once \(h-1\ge n\).

Writing \(M=2^h\), the number of empirical types is at most

\[
 \binom{n+M-1}{M-1}\le[e(1+n/M)]^M.
 \tag{3.2}\label{eq:3.2}
\]

Indeed, the nonnegative block counts sum to \(n\). Stars and bars gives
the first bound in \eqref{eq:3.2}; then
\(\binom{n+M-1}{M-1}\le\binom{n+M}{M}\le[e(n+M)/M]^M\). No assumption
that different component types are independent is used.

\subsection{Mixture entropy}\label{mixture-entropy}

For a valid family, sample a position uniformly among its \(L\)
positions and read the resulting stationary binary process. All
\(d\)-blocks are distinct, so

\[
 H_d=\log_2L,\qquad H_d-H_{d-1}=\frac{2m}{L}.
\]

Stationary conditional entropies decrease with block length. Thus

\[
 \delta_h\le\frac{H_h}{h}
 \le\frac{\log_2L-2m(d-h)/L}{h}.
 \tag{3.3}\label{eq:3.3}
\]

Let \(J\) identify the component circle and let \(X,Y\) be the
preceding \(h\)-1 bits and the final bit. Then

\[
 L\delta_h^{\rm mix}-\sum_i r_i\delta_h(w_i)
 =L I(Y;J\mid X)\ge0.
 \tag{3.4}\label{eq:3.4}
\]

This remains valid for any number of component circles and arbitrary
lengths. Applying \eqref{eq:3.1}--(3.4), and bounding the rooted prefactor by
\(L^k\), yields

\[
 \log_2Q_d^{(m)}(\mathbf r)
 \le \frac{Lt}{h}
 +kM\log_2[e(1+L/M)]+k\log_2L-\kappa m,
 \tag{3.5}\label{eq:3.5}
\]

where

\[
 t=\log_2L-h,\qquad \kappa=\frac{2d-3h}{h}.
\]

\subsection{The sparse threshold}\label{the-sparse-threshold}

Suppose \(d^4\le L\le N^\beta\), and put \(q=\log_2d\) and

\[
 h=\lfloor\log_2L-\log_2d-\log_2q\rfloor.
\]

For sufficiently large \(d\), \(h\ge2\) and \(h\le\beta d\),

\[
 \kappa\ge(2-3\beta)/\beta>0,
 \qquad kM\log_2[e(1+L/M)]\le 2kL/d.
\]

Let

\[
 F=Lt/h+2kL/d+k\log_2L,\qquad
 m_0=\left\lceil\frac{F+L/d}{\kappa}\right\rceil.
\]

Summing \eqref{eq:3.5} gives

\[
 D_{\mathbf r}:=\sum_{m\ge m_0}Q_d^{(m)}(\mathbf r)
 \le C_\beta2^{-L/d}.
 \tag{3.6}\label{eq:3.6}
\]

Furthermore,

\[
 \frac{2dm_0}{L}
 \le p_\beta(\log_2d+\log_2\log_2d)+O_{k,\beta}(1).
 \tag{3.7}\label{eq:3.7}
\]

For clarity, the algebra behind (3.7) is \[
 \frac{2dm_0}{L}
 \le \frac{2d}{2d-3h}
 \left[t+(2k+1)\frac hd+\frac{kh\log_2L}{L}\right]
 +\frac{2d}{L}.
\] The leading multiplier is at most \(p_\beta\);
\(t\in[\log_2d+\log_2\log_2d,\,
\log_2d+\log_2\log_2d+1)\), \(h/d\le\beta\), and
\(h\log_2L/L=O_\beta(d^{-2})\) because \(L\ge d^4\). Also, summation in
(3.6) uses the finite positive lower bound on \(\kappa\):
\(\sum_{m\ge m_0}2^{F-\kappa m}\le
2^{-L/d}/(1-2^{-\kappa})\).

When \(r=\max_i r_i\), \(L/r\le k\). Hence the closing penalty at the
longest cycle's scale satisfies

\[
 2^{2dm_0/r}\le C_{k,\beta}(d\log_2d)^{kp_\beta}.
 \tag{3.8}\label{eq:3.8}
\]

For \(L<d^4\) no dense decomposition is necessary. Stationarity also
gives

\[
 m\le\frac{L\log_2L}{2d},
\]

so, for any \(m_{\rm new}\le m\),

\[
 2^{2dm_{\rm new}/r}\le L^{L/r}\le L^k\le d^{4k}.
 \tag{3.9}\label{eq:3.9}
\]

In either case the needed sparse penalty is bounded by

\[
 H_{k,\beta}(d)=C_{k,\beta}\left[d^{4k}+(d\log_2d)^{kp_\beta}\right].
 \tag{3.10}\label{eq:3.10}
\]

It depends on \(k\), not on the ratios between the component lengths.

\section{Marked stopped collisions with earlier cycles
present}\label{marked-stopped-collisions-with-earlier-cycles-present}

Let \(\mathbf a=(r_1,\ldots,r_{k-1})\), \(P=\sum_{i=1}^{k-1}r_i\), and let
\(\mu_{\mathbf a}\) be its scaled ordinary cycle measure. Conditional on
the recorded rooted cycles, feedback values at all unqueried contexts
remain independent fair bits: specifying the recorded cycles imposes
exactly the feedback constraints along them.

Adjoin a new independent uniform root \(X\) and run its ordinary
trajectory for \(n\) steps. Retain it only if its first \(n+1\) states
are distinct and outside all the earlier cycles. Let \(J_{\rm int}\)
count repeated contexts inside that prefix and \(J_{\rm cross}\) count
prefix queries using contexts of earlier cycles. The latter statistic
may be bounded after dropping the retention restriction.

\subsection{Internal contacts and marked
switching}\label{internal-contacts-and-marked-switching}

For \(h>0\), the conjugate endpoint count must include the condition
\(L_f(y)>h\). On an internal conjugate encounter of the new surviving
path, both states of the conjugate pair lie in a cycle disjoint from all
earlier cycles. Their context is therefore unused by every earlier
cycle. Flipping its feedback bit preserves all earlier marked cycles and
splits off an \(h\)-cycle.

The resulting marked configurations are a subset of those counted by
\(Q_d(\mathbf a,h)\). Stationarity of the new uniform root then gives

\[
 \mathbb E_{\mu_{\mathbf a},X}J_{\rm int}
 \le\frac1N\sum_{h=1}^{n-1}(n-h)Q_d(\mathbf a,h).
 \tag{4.1}\label{eq:4.1}
\]

This is not a conditional-uniformity assertion about \(f\) under
\(\mu_{\mathbf a}\). It follows from the feedback-flip bijection while
keeping the previously selected cycles marked. The survival and
disjointness restrictions ensure the flip leaves those marks unchanged.
Sorting the list \(({\mathbf a},h)\) is permitted because \(Q\) is
symmetric.

Here is the normalization and bijection in \eqref{eq:4.1}. Write
\(\mathcal C_f(y)\) for the cycle containing \(y\) and \(A\) for the
union of the previously marked cycles. Feedback-bit switching gives \[
 \int\!\mu_{\mathbf a}(d(f,\text{marks}))
 \sum_{\substack{y:T_f^hy=\bar y,\ L_f(y)>h\\
                  \mathcal C_f(y)\cap A=\varnothing}}1
 \le Q_d(\mathbf a,h).
 \tag{4.1a}\label{eq:4.1a}
\] To see this, flip the bit at \(u(y)\). In the new permutation the
cycle containing \(\bar y\) has length \(h\), and \(y\) is outside it.
All earlier marked cycles, including their roots, are unchanged. The map
is injective: from the new marked \(h\)-cycle rooted at \(\bar y\),
recover \(y\), flip the same bit, and recover the original
configuration. Counting all marked \(h\)-cycles can only enlarge the
image. The uniform feedback law is invariant under this flip. The factor
\(h\) in \(Q_d(\mathbf a,h)\) counts its possible new roots, including
when \(h\) equals one of the earlier lengths.

For a retained prefix and query times \(i<j\), take \(y=T_f^iX\) and
\(h=j-i\). Conditional on \(f\) and the earlier marks, \(y\) is uniform.
The retained prefix is in a cycle disjoint from \(A\); disjointness of a
prefix from \(A\) implies this because \(A\) is a union of complete
cycles. It also has \(L_f(y)>n>h\). Dropping the stronger retention
conditions and applying \eqref{eq:4.1a} gives \(Q_d(\mathbf a,h)/N\) for each
pair. There are \(n-h\) query pairs with separation \(h\), proving
(4.1). No unconditioned stationarity assertion is substituted for this
marked calculation.

\subsection{Cross contacts}\label{cross-contacts}

For fixed \(f\) and fixed earlier cycles, every iterate of \(X\) is
uniform because \(T_f\) is a permutation. The earlier cycles use at most
\(P\) contexts. Each context has exactly two possible preceding states,
so

\[
 \mathbb E_{\mu_{\mathbf a},X}J_{\rm cross}
 \le\frac{2nP}{N}Q_d(\mathbf a).
 \tag{4.2}\label{eq:4.2}
\]

On a retained prefix, each context occurs at most twice across its query
states and the earlier cycles: a context has only two preceding states,
and all these states are distinct. Hence a prefix query cannot
simultaneously repeat an earlier prefix query and a context in the
marked cycles. Thus the reused-query count is exactly
\(J=J_{\rm int}+J_{\rm cross}\) on retained histories. Both quantities
are set to zero on discarded histories; only the cross bound is enlarged
by dropping retention.

Consequently the ordinary open-repeat statistic \(J\) satisfies

\[
 \mathbb E J\le
 \frac1N\sum_{h=1}^{n-1}(n-h)Q_d(\mathbf a,h)
 +\frac{2nP}{N}Q_d(\mathbf a).
 \tag{4.3}\label{eq:4.3}
\]

The first term involves the same number of cycles but strictly smaller
total length. The second involves one fewer cycle. This is the induction
structure.

\section{Longest-last forced closure}\label{longest-last-forced-closure}

Sort \(r_1\le\cdots\le r_k=r\). Keep the earlier cycles under
\(\mu_{\mathbf a}\). Starting from a new uniform root, expose \(r-d\)
ordinary open steps, then force the final \(d\) outputs to equal the
initial state. Never change an already prescribed feedback value.
Success requires first return at exactly \(r\) and disjointness from all
earlier cycles. Complete unqueried feedback values fairly.

Operationally, first sample the \emph{recorded} marked-cycle histories
from the projection of \(\mu_{\mathbf a}\), which is a finite measure.
Sample only the still-unqueried feedback values as the new trajectory
needs them, and complete the remaining values at the end. This avoids
treating the biased law after an earlier forced closure as an ordinary
feedback law. Every previously selected cycle constrains only its own
queried contexts.

Denote this finite successful measure by \(\nu_{\mathbf a;r}\). Its mass
is at most \(Q_d(\mathbf a)\), not necessarily at most one. Let \(K\)
count closing queries whose feedback value has already been fixed
strictly before that query is made: by an earlier marked cycle, by the
current open path, or by an earlier query in the current closing
segment. In particular, two occurrences of a context within the closing
segment contribute one to \(K\). On successful histories,

\[
 d\nu_{\mathbf a;r}=2^{-K}d\mu_{\mathbf a,r},
 \tag{5.1}\label{eq:5.1}
\]

For an atom with \(v_{\rm old}\) old contexts and \(v\) contexts in the
completed family, the scaled ordinary word weight is \(2^{-v}\). The
forced construction samples the new root with factor \(N^{-1}\) and
samples only \(v-v_{\rm old}-(d-K)\) additional feedback bits fairly.
Its weight is therefore \(2^{-v_{\rm old}}N^{-1}
2^{-[v-v_{\rm old}-(d-K)]}=2^{-v-K}\). Unused fair completion is
identical. This proves \eqref{eq:5.1} at the atom level.

For a fixed full cycle family, let \(m_{\rm new}\) count all repeats
involving the last cycle: internal repeats plus contexts it shares with
earlier cycles. Earlier cycles' internal repeats are not counted in
\(m_{\rm new}\). If \(J\) counts reused queries in the last cycle's open
part (with earlier contexts known at time zero), then

\[
 J+K=m_{\rm new}.
 \tag{5.2}\label{eq:5.2}
\]

Average only the root of the last cycle. A shared context belongs to its
closing part with probability \(d/r\); an internal repeated context
contributes with probability at most \(2d/r\). Thus

\[
 \mathbb E_{\rm root}K\le2dm_{\rm new}/r,
 \quad \mathbb E_{\rm root}J\ge a_rm_{\rm new},
 \quad a_r=1-2d/r.
\]

The function \(j\mapsto j2^{j-m_{\rm new}}\) is increasing and convex
for \(j\ge0\). Therefore

\[
 \mathbb E_{\rm root}[J2^{-K}]
 \ge a_rm_{\rm new}2^{-2dm_{\rm new}/r}.
 \tag{5.3}\label{eq:5.3}
\]

On the sparse class, (3.8) or (3.9) bounds the inverse factor by
\(H_{k,\beta}\). The forced open part has exactly the ordinary
\(\mu_{\mathbf a}\)-plus-uniform-root law; success only restricts it. To
make the integration explicit, let \(A_{\rm sp}\) be the sparse class,
\(\mu=\mu_{\mathbf a,r}\), \(\nu=\nu_{\mathbf a;r}\), and
\(H=H_{k,\beta}(d)\). Then \[
 \begin{aligned}
 \frac{a_r}{H}\int_{A_{\rm sp}}m_{\rm new}\,d\mu
 &\le \int_{A_{\rm sp}}J2^{-K}\,d\mu\\
 &=\int_{A_{\rm sp}}J\,d\nu
 \le \mathbb E_{\mu_{\mathbf a},X}J.
 \end{aligned}
 \tag{5.3a}\label{eq:5.3a}
\] The sparse class and \(m_{\rm new}\) are invariant under re-rooting
the last cycle, so the first inequality follows by integrating the
root-averaged bound \eqref{eq:5.3}. The last inequality discards a restriction to
successful histories; it does not assert that the feedback law after
forced closure is ordinary. Using \(1-2^{-K}\le K\) now gives

\[
\begin{split}
0\le Q_d(\mathbf a,r)-\nu_{\mathbf a;r}(\Omega)
&\le\frac{2d}{r}\int_{\rm sparse}m_{\rm new}\,d\mu_{\mathbf a,r}+D_{\mathbf a,r}\\
&\le\frac{2dH_{k,\beta}}{r-2d}\mathbb E J+D_{\mathbf a,r}.
\end{split}
\tag{5.4}\label{eq:5.4}
\]

Take \(r\ge4d\) and \(n=r-d\). Combining \eqref{eq:4.3}, (5.4), and
\(\nu_{\mathbf a;r}(\Omega)\le Q_d(\mathbf a)\), yields the key finite
recurrence

\[
\boxed{
Q_d(\mathbf a,r)\le Q_d(\mathbf a)
+\frac{2dH_{k,\beta}}{r-2d}
 \left[
 \frac1N\sum_{h=1}^{n-1}(n-h)Q_d(\mathbf a,h)
 +\frac{2nP}{N}Q_d(\mathbf a)
 \right]+D_{\mathbf a,r}.
}
\tag{5.5}\label{eq:5.5}
\]

\(D=0\) for total length below \(d^4\); otherwise it is bounded by
(3.6).

No long cycle is charged at a shorter cycle's rate. Prior cycles remain
under their ordinary finite measure, rather than being incorrectly
treated as an unbiased probability law after earlier forced closures.

\section{Close the upper induction}\label{close-the-upper-induction}

\subsection{\texorpdfstring{Base: all lengths below
\(4d\)}{Base: all lengths below 4d}}\label{base-all-lengths-below-4d}

Let \(B\) mean that every prescribed cyclic word is primitive and no two
words represent the same necklace. \(A_d\) implies \(B\). On \(B\) with
no repeated contexts, the weighted integrand is exactly one.

For two distinct positions, Appendix A proves the fully wrapped-window
bound \[
 \mathbb P\{B\text{ and their }(d-1)\text{-contexts agree}\}
 \le 2^{-(d-1)}.
\] This includes component words shorter than the window and equal
component lengths; the event \(B\) removes the periodic degeneracies.

A union bound, together with \(m\le L\log_2L/(2d)\), implies

\[
 Q_d(\mathbf r)\le1+\frac{L^{2k+2}}N,
 \qquad \max_i r_i<4d,\quad L<4kd.
 \tag{6.1}\label{eq:6.1}
\]

\subsection{Crude bound}\label{crude-bound}

Induct first on the number of cycles \(k\), then strongly on total
length \(L\) for fixed \(k\). The target crude bound is
\(Q_d(\mathbf r)\le 2^k\). The base (6.1) holds for large \(d\). When
\(r=\max_i r_i\ge4d\), all same-\(k\) moments in \eqref{eq:5.5} have total length
\(P+h<L\). The lower-order moment is bounded by \(2^{k-1}\). Therefore

\[
 Q_d(\mathbf a,r)\le2^{k-1}
 +C_k dH_{k,\beta}(d)\frac LN+D_{\mathbf a,r}.
\]

The extra terms tend to zero uniformly for \(L\le N^\beta\), since \(H\)
is polynomial and \(\beta<1\). Thus they fit within the fixed gap from
\(2^{k-1}\) to \(2^k\). This closes the induction, including moment
lists that contain very short components. In particular, \(Q_d(h)\le2\)
for all \(h\le N^\beta\).

This is a strong induction on an integer total length, within an outer
induction on the fixed order. For a fixed target order \(k\), take \(d\)
large enough for all orders at most \(k\). At each non-base step the
strictly smaller-order term is bounded by \(2^{k-1}\), while the
same-order terms have strictly smaller total. Thus there is a fixed gap
available for the uniformly vanishing error; no unproved sharp estimate
is used to initialize the recurrence.

\subsection{Sharp constant}\label{sharp-constant}

Substitute the crude bounds back into \eqref{eq:5.5}:

\[
 Q_d(\mathbf a,r)\le Q_d(\mathbf a)+C_{k,\beta}dH_{k,\beta}(d)\frac LN+D_{\mathbf a,r}.
\]

Repeatedly remove a longest component. There are at most \(k\) steps.
The base (6.1) costs at most a polynomial times \(1/N\), absorbed by
\(P_{k,\beta}(d)L/N\). For \(L\ge d^4\) the dense error is at most
\(C_\beta2^{-d^3}\), also absorbed. This proves \eqref{eq:1.1}.

\section{A matching lower bound from several middle-first
bridges}\label{a-matching-lower-bound-from-several-middle-first-bridges}

Only the one-cycle bound \(Q_d(h)\le2\) is needed in this section. The
restriction \(L\le N^\beta\) is used here only to ensure this bound at
the required lengths. Thus the numerical lower-bound constant depends
only on \(k\), although the threshold for sufficiently large \(d\) may
also depend on \(\beta\).

For an ordinary trajectory of \(n\ge d\) steps from a uniform root, let
\(R_{n,d}\) mean that some query in its final \(d\) steps repeats an
earlier context. The stopped conjugate-switch identity and a union bound
over time separations give

\[
 \mathbb P(R_{n,d})\le\frac{d+1}{N}\sum_{h\le n}Q_d(h)
 \le\frac{2(d+1)n}{N}.
 \tag{7.1}\label{eq:7.1}
\]

The event of a prior return costs \(N^{-1}\sum_{h\le n}Q_d(h)\); on
survival a repeated context is a conjugate encounter, and each
separation appears at most \(d\) times in the terminal query block.
Coupling the last \(d\) fresh outputs to fair bits also gives

\[
 d_{\rm TV}(\mathcal L(X_0,X_n),U_N\otimes U_N)
 \le\frac{2(d+1)n}{N}.
 \tag{7.2}\label{eq:7.2}
\]

Coordinate reversal conjugates the inverse Thorp permutation to the same
model with its feedback arguments reversed, so these estimates hold
backwards too.

More explicitly, the one-cycle instance of \eqref{eq:4.1a} bounds the stopped
conjugate endpoint count by \(Q_d(h)\). With a uniform root, a prior
return has probability \(N^{-1}\sum_{h\le n}Q_d(h)\). On survival, a
pair of equal contexts consists of distinct conjugate states. For each
separation \(h\), there are at most \(d\) pairs with the later query
among the final \(d\) queries. This proves \eqref{eq:7.1}.

For (7.2), at each of the final \(d\) steps supply a new independent
fair comparison bit. If its context is new, choose the freshly revealed
feedback value to make the actual output equal that bit. If the context
is old, continue the comparison with its independent bit. The comparison
endpoint is a uniform \(d\)-bit state independent of the initial root;
failure of the coupling requires terminal reuse. This constructs the
comparison, rather than conditioning a collection of outputs on the
event that all are fresh.

If \(R_d\) reverses all \(d\) coordinates, direct calculation gives \[
 R_dT_f^{-1}R_d=T_{f\circ R_{d-1}}.
\] The transformed feedback table is again uniform and independent, so
the same stopped-reuse and endpoint estimates hold backwards.

\subsection{Reference construction}\label{reference-construction}

Assume all \(r_i\ge4d\) and put \(t_i=r_i-2d\). In a common ordinary
feedback environment, choose independent uniform roots \(v_i\) and run
middle paths of \(t_i\) steps to \(z_i\). Independently choose uniform
\(w_i\). Add the two unique \(d\)-step shift paths

\[
 z_i\longrightarrow w_i\longrightarrow v_i.
\]

Let \(G\) require all middle states to be distinct, within and between
paths, and all \(2kd\) added query contexts to be mutually distinct and
absent from every middle queried context. On \(G\) assign the fresh
feedback values to realize the added outputs and complete all remaining
values fairly. This produces \(k\) disjoint cycles of the required
lengths.

\subsection{Failure bounds}\label{failure-bounds}

Internal failure of a middle path costs at most \(2t_i/N\).

For different middle paths of lengths \(t_i,t_j\), conditional on any
fixed \(f\), an intersection entails

\[
 v_j=T_f^{a-b}v_i,
 \qquad 0\le a\le t_i,\quad0\le b\le t_j.
\]

There are at most \(t_i+t_j+1\) distinct offsets. Thus

\[
 \mathbb P(\text{middle }i\text{ meets middle }j\mid f)
 \le(t_i+t_j+1)/N.
 \tag{7.3}\label{eq:7.3}
\]

There is no \(t_it_j/N\) loss here.

For an added path, separate contacts with its \textbf{own} middle from
contacts with \textbf{other} middles. For self-contacts and contacts
with its own middle, consider just that single ordinary middle and its
independent seed \(w_i\). Until the first reused query, its forward
reference extension agrees with an ordinary extension of that middle, by
deferred decisions. Formula \eqref{eq:7.1} with \(n=t_i+d\) bounds this failure.
The other extension is explored backwards from \(v_i\), and the
reverse-time version of \eqref{eq:7.1} gives the same bound. A backward step from
\(y=(u,z)\) has predecessor \((z\oplus f(u),u)\) and queries exactly
\(f(u)\), the same entry as the forward step into \(y\). Hence the
reverse-time reuse event controls precisely the relevant forward-context
contacts. These are marginal, unconditional estimates; they do not
require the feedback law to remain ordinary after conditioning on every
other trace.

For contacts with a different middle \(j\), use the reference
construction directly. Conditional on \(f\), every fixed context in the
added path \(z_i\) to \(w_i\) is uniform: it uses a suffix of the
uniform \(z_i\) and a prefix of the independent uniform \(w_i\). The
path \(w_i\) to \(v_i\) has the same property. These contexts are
independent of a specified query context in middle \(j\), since
\((v_i,w_i)\) and \(v_j\) are independent conditional on \(f\). Equality
therefore has probability \(2/N\). A union bound over at most \(d t_j\)
pairs for each added path costs \(2dt_j/N\). Summing all the self and
other-middle bounds costs \(O_k(dL/N)\).

For the two added paths belonging to the same cycle, their possible
cross contacts depend only on \((z_i,w_i,v_i)\). By (7.2) this
triple is within \(2(d+1)t_i/N\) of three independent uniform states.
Under the independent law, \(z_iw_iv_i\) is a linear independent-bit
word; each of the \(d^2\) possible cross-window matches has probability
\(2/N\), including overlapping windows. This again sums to
\(O_k(dL/N)\).

For added paths belonging to different cycles, no endpoint approximation
is needed. Conditional on \(f\), the pairs \((v_i,w_i)\) are independent
across \(i\), \(z_i\) is uniform, and each individual added context is
uniform. Thus each cross-cycle pair of added contexts agrees with
probability \(2/N\). There are \(O(k^2d^2)\) such pairs, costing
\(O(k^2d^2/N)\), absorbed into \(O_k(dL/N)\).

All estimates are unconditional union bounds; the failure events are not
assumed independent. Together,

\[
 \mathbb P_{\rm ref}(G)\ge1-C_k dL/N.
 \tag{7.4}\label{eq:7.4}
\]

\subsection{Exact good-history weight}\label{exact-good-history-weight}

Fix the middle traces, whose queried contexts number \(v_M\), and the
words \(w_i\). The ordinary root-and-middle probability contains
\(N^{-k}2^{-v_M}\). On \(G\) the \(2kd\) prescribed fresh outputs have
probability \(N^{-2k}\); scaling the ordinary law by \(N^k\) leaves
weight \(2^{-(v_M+2kd)}\). The reference chooses all \(w_i\) with
probability \(N^{-k}\), giving the same weight \(N^{-2k}2^{-v_M}\).

The reference construction only retains feedback values queried by the
middles. Values at added contexts are assigned after choosing the
bridges, and the rest are freshly completed. It does not require an
initially sampled unused feedback value to equal a forced value; those
initially unused values have been marginalized. This distinction is part
of the weight calculation, not an assumption about their conditional
distribution.

Thus the scaled ordinary cycle measure equals the reference measure on
\(G\), including the same uniform completion of unused feedback values.
Therefore

\[
 Q_d(\mathbf r)\ge\mathbb P_{\rm ref}(G)\ge1-C_kdL/N.
\]

This proves the lower half of \eqref{eq:1.2}. It needs neither conditional
independence of endpoints given all other traces nor a joint-collision
moment theorem.

\section{Fixed-order moments and the truncated Gaussian
field}\label{fixed-order-moments-and-the-truncated-gaussian-field}

Fix \(0<\alpha<2/3\) and choose \(\beta\) strictly between \(\alpha\)
and \(2/3\). Let \(b=\lfloor N^\alpha\rfloor\) and \(\ell=\log N\). For
every fixed moment order \(m\) and all sufficiently large \(d\), all
lists of at most \(m\) lengths bounded by \(b\) have total length
\(\le N^\beta\). The moment comparison in this section is initially
restricted to \(4d\le r\le b\); the omitted shorter lengths are restored
after the tightness argument. On this restricted range, (1.2) gives

\[
 \max_{1\le j\le m}\ \sup_{4d\le r_1,\ldots,r_j\le b}
 |Q_d(r_1,\ldots,r_j)-1|
 \le \varepsilon_{d,m},\qquad
 \varepsilon_{d,m}=d^{O_{m,\beta}(1)}N^{\alpha-1}.
 \tag{8.1}\label{eq:8.1}
\]

This is smaller than every negative power of \(d\) for fixed \(m\).

Let \(Z_r\), \(4d\le r\le b\), be independent Poisson variables of means
\(1/r\). Their normalized factorial moments \(Q\) are exactly one.
Ordinary moments of additive cycle statistics expand over partitions of
their selected cycle occurrences into groups that refer to the same
cycle. Each block contributes one harmonic summation. For logarithmic
weights, a block of multiplicity a contributes at most

\[
 \sum_{r\le b}(\log r)^a/r=O_a(\ell^{a+1}).
\]

Consequently the error in a fixed raw \(m\)-th moment of the logarithmic
product is \(O_m(\varepsilon_{d,m}\ell^{2m})\). Centering and Gaussian
normalization multiply this only by further powers of \(\ell\). All
normalized moment errors vanish. The same argument applies to arbitrary
finite collections of count and log-product increments. The
corresponding Poisson statistics have Gaussian limits and
moment-determinate limiting distributions.

For completeness, the reference Gaussian limit itself needs no
permutation-order theorem. At positions \(t_r=\log r/\ell\), a linear
combination of finitely many count and normalized log-product increments
has centered Poisson form \[
 \frac1{\sqrt\ell}\sum_{4d\le r\le b}
        g(t_r)(Z_r-1/r),
\] where \(g\) is a bounded, piecewise continuous function. Its variance
converges to \(\int_0^\alpha g(t)^2\,dt\) by harmonic Riemann sums. Its
cumulant of order \(j\ge3\) is \[
 \ell^{-j/2}\sum_{4d\le r\le b}\frac{g(t_r)^j}{r}
 =O_j(\ell^{1-j/2})=o(1).
\] Thus its moments converge to Gaussian moments. The Gaussian law is
moment-determinate; the multivariate statement follows by the
Cramér--Wold argument. The log-product coordinate corresponds to the
weight \(\sqrt3\,t\), so the common limiting Brownian motion gives
\(\bigl(W(u),\sqrt3\int_0^u t\,dW(t)\bigr)\).

\begin{corollary}[Bounded additive tests]\label{cor:additive}
Fix $0<\alpha<2/3$, put $\ell=\log N$, and let
$g_1,\ldots,g_q$ be a fixed finite family of bounded piecewise-continuous
real functions on $[0,\alpha]$. Then
\[
 \left(
 \frac1{\sqrt\ell}\sum_{4d\le r\le N^\alpha}
 g_j\!\left(\frac{\log r}{\ell}\right)(C_r-1/r)
 \right)_{j=1}^q
 \Longrightarrow \mathcal N_q(0,\Sigma),
 \qquad
 \Sigma_{ij}=\int_0^\alpha g_i(t)g_j(t)\,dt.
\]
\end{corollary}
\begin{proof}
Apply the preceding factorial-moment comparison to each linear combination
of the $g_j$. The Poisson reference has variance converging to the stated
integral and cumulants of every fixed order $m\ge3$ of size
$O_m(\ell^{1-m/2})$. Its normalized moments converge to Gaussian moments;
the comparison transfers them. Cram\'{e}r--Wold gives the joint conclusion,
including singular covariance matrices.
\end{proof}


\subsection{Tightness, not just finite-dimensional
convergence}\label{tightness-not-just-finite-dimensional-convergence}

First omit lengths below \(4d\). Center by the exact Poisson-reference
harmonic sums. In the log-product coordinate use weights
\(\sqrt3\log r/\ell\), so both coordinates are sums with bounded
weights, divided by \(\sqrt\ell\).

Extend each truncated centered process constantly beyond \(\alpha\) to
the next point of the full grid \(u=j/\ell\). Interpolate on this full
grid and then restrict to \([0,\alpha]\). A grid interval containing
\(m\ge1\) full mesh steps has reference Poisson intensity \(O(m)\). For
either normalized coordinate its fourth centered Poisson moment is at
most \[
 \frac{C(m^2+m)}{\ell^2}.
\] Expansion into raw moments of orders at most four and \eqref{eq:8.1} gives the
same bound for the actual centered coordinate, plus a uniform remainder
\(R_d\) that is smaller than every power of \(1/\ell\). For example, the
crude bound \(R_d=O(\varepsilon_{d,4}(1+\ell)^6)\) suffices. Since
\(m\ge1\), absorb this remainder into a constant multiple of
\((m/\ell)^2\).

Linearly interpolate the grid values. Within a grid cell an increment is
a fraction \(\theta\) of the grid increment. Its fourth moment is at
most \(C\theta^4/\ell^2\le C|u-v|^2\), since \(\theta=\ell|u-v|\le1\).
Across cells split an interval into its first partial cell, full grid
part, and final partial cell. The inequality
\(|x+y+z|^4\le27(|x|^4+|y|^4+|z|^4)\) gives the uniform continuity bound
\(C|u-v|^2\). The constant extension makes the final mesh cell a full
cell as well; no uniform estimate for an arbitrarily shortened
interpolation cell is needed. The fourth-moment tightness criterion
therefore applies to the interpolated processes.

To compare with the original step processes, let \(D_j\) count all
cycles in one logarithmic grid bin. Its harmonic intensity is bounded by
an absolute constant. The upper moment theorem at orders at most four
gives \(\mathbb ED_j^4\le C\), uniformly over the bins. There are
\(O(\ell)\) bins, so for every \(\eta>0\), \[
 \mathbb P\{\max_jD_j>\eta\sqrt\ell\}
 \le \frac{C}{\eta^4\ell}\longrightarrow0.
\] Within a bin, the normalized log-product oscillation is at most a
fixed multiple of \(D_j/\sqrt\ell\), because \(\log r/\ell\le\alpha\).
The deterministic reference centerings have oscillation
\(O(1/\sqrt\ell)\) there. Thus both step processes are uniformly close
in probability to the tight interpolations.

Below length \(4d\), \(Q_d(r)\le2\) implies expected count \(O(\log d)\)
and expected log-product \(O((\log d)^2)\). Both vanish at their
respective normalizations. This also handles the initial interval near
\(u=0\). Centering by exact harmonic sums or by the displayed polynomial
centerings differs by a uniformly negligible quantity.

Therefore, with

\[
 K_d(u)=\sum_{r\le N^u}C_r,\qquad
 B_d(u)=\sum_{r\le N^u}C_r\log r,
\]

we obtain convergence in \(D([0,\alpha],\mathbb R^2)\), with continuous
limiting paths:

\[
 \left(
 \frac{K_d(u)-u\ell}{\sqrt\ell},\quad
 \frac{B_d(u)-\tfrac12u^2\ell^2}{\sqrt{\ell^3/3}}
 \right)
 \Longrightarrow
 \left(W(u),\sqrt3\int_0^u t\,dW(t)\right).
 \tag{8.2}\label{eq:8.2}
\]

This argument does not require a total-variation approximation of the
full initial cycle-count vector. Fixed-order estimates with
superpolynomially small errors suffice for the moment and tightness
argument. Section 10 separately obtains total-variation approximation on
sets of bounded harmonic mass.

\section{Product to LCM}\label{product-to-lcm}

For all \(r,s\le b\), the all-positive-length upper theorem gives, for
large \(d\),

\[
 \mathbb EC_r\le2/r,\qquad
 \mathbb E[C_r(C_s-\mathbf1_{r=s})]\le2/(rs).
\]

Let \(D_m=\sum_{m\mid r,\ r\le b}C_r\). The exact arithmetic identity is

\[
 B_d(\alpha)-\log\operatorname{lcm}\{r\le b:C_r>0\}
 =\sum_{m\le b}\Lambda(m)(D_m-1)_+.
\]

The preceding harmonic bounds imply

\[
 \mathbb E(D_m-1)_+
 \le\min\left\{\frac{2H_{\lfloor b/m\rfloor}}m,
 \frac{H_{\lfloor b/m\rfloor}^2}{m^2}\right\}.
\]

Put \(t=\log b\). The elementary Chebyshev bound
\(\psi(x)=\sum_{m\le x}\Lambda(m)=O(x)\) and partial summation give \[
 \sum_{m\le t}\frac{\Lambda(m)}m=O(\log t),\qquad
 \sum_{m>t}\frac{\Lambda(m)}{m^2}=O(1/t).
 \tag{9.1}\label{eq:9.1}
\] For completeness, the product of primes \(p\) for which some power
satisfies \(n<p^a\le2n\) divides \(\binom{2n}{n}\), and there is at most
one such power for each prime. Thus
\(\psi(2n)-\psi(n)\le\log\binom{2n}{n}\le2n\log2\); summing over dyadic
intervals proves \(\psi(x)=O(x)\). In the harmonic bound above,
\(H_{\lfloor b/m\rfloor}=O(t)\) uniformly. Splitting at \(m=t\)
therefore yields \[
 \begin{aligned}
 \mathbb E\bigl(B_d(\alpha)-O_d(\alpha)\bigr)
 &\ll t\sum_{m\le t}\frac{\Lambda(m)}m
       +t^2\sum_{m>t}\frac{\Lambda(m)}{m^2}\\
 &=O(\log b\,\log\log b)=o(\ell^{3/2}).
 \end{aligned}
 \tag{9.2}\label{eq:9.2}
\]

The discrepancy is nondecreasing with the cutoff, so the bound is
uniform on \(0\le u\le\alpha\). Markov's inequality transfers (8.2) from
\(B_d(u)\) to

\[
 O_d(u)=\log\operatorname{lcm}\{r\le N^u:C_r>0\}.
\]

This proves Theorem~\ref{thm:functional}, including its scalar consequence.

This is a truncated theorem. The dense-tail estimate used above requires
\(\kappa>0\) in \eqref{eq:3.5}; the argument does not establish the full
permutation-order limit.


\begin{corollary}[Lower bounds for the full permutation]\label{cor:full-lower}
For every $\varepsilon>0$,
\[
 \mathbb P\!\left\{
 \log\operatorname{ord}(T_f)
 \ge(2/9-\varepsilon)(\log N)^2\right\}\longrightarrow1,
 \qquad
 \mathbb P\!\left\{\sum_{r=1}^N C_r
 \ge(2/3-\varepsilon)\log N\right\}\longrightarrow1.
\]
\end{corollary}
\begin{proof}
For each fixed $\alpha<2/3$, Theorem~\ref{thm:functional} gives the laws of
large numbers $O_d(\alpha)/\ell^2\to\alpha^2/2$ and
$K_d(\alpha)/\ell\to\alpha$ in probability. The full order and count
dominate their truncated counterparts. Choose a fixed $\alpha$ sufficiently
close to $2/3$. No uniform estimate as $\alpha\uparrow2/3$ is used.
\end{proof}


\section{Poisson limits and further cycle laws}\label{sec:poisson}
The moment theorem gives a growing-dimensional total-variation result when
the total harmonic mass of the observed lengths is bounded. This is
different from approximating the entire initial vector through a power of $N$.

\begin{theorem}[Bounded harmonic mass]\label{thm:bounded-mass}
Fix $0<\alpha<2/3$ and $\Lambda<\infty$. Let
$I_d\subseteq\{r:4d\le r\le N^\alpha\}$ satisfy
$\lambda_d=\sum_{r\in I_d}r^{-1}\le\Lambda$. Then
\[
 d_{\rm TV}\!\left(
 \mathcal L((C_r)_{r\in I_d}),
 \bigotimes_{r\in I_d}\operatorname{Poisson}(1/r)
 \right)\longrightarrow0.
 \tag{10.1}\label{eq:10.1}
\]
The cardinality of $I_d$ may tend to infinity.
\end{theorem}

\subsection{Finite inclusion--exclusion}
We give a finite comparison that makes the order of limits explicit.
For a vector $X=(X_i)_{i\in I}$ of nonnegative integers on a finite set,
write its ordered factorial moment measures as
\[
 \mu_j(i_1,\ldots,i_j)
 =\mathbb E\prod_{i\in I}(X_i)_{a_i},
 \qquad a_i=\#\{v:i_v=i\}.
\]
Let $Y_i$ be independent Poisson variables with means $\lambda_i$, put
$\lambda=\sum_i\lambda_i$, and let
$\pi_j(i_1,\ldots,i_j)=\prod_v\lambda_{i_v}$.
The norm $\|\cdot\|_1$ below is total absolute mass on $I^j$.

\begin{lemma}[Finite factorial-moment comparison]\label{lem:finite-ie}
If $\mu_p(I^p)<\infty$ for an integer $p\ge1$, then
\[
 \begin{aligned}
 d_{\rm TV}(\mathcal L(X),\mathcal L(Y))
 \le{}&\sum_{j=1}^{p-1}\frac{2^j}{j!}\|\mu_j-\pi_j\|_1\\
 &+\frac{2^p}{p!}\bigl(\mu_p(I^p)+\lambda^p\bigr).
 \end{aligned}
 \tag{10.2}\label{eq:10.2}
\]
\end{lemma}
\begin{proof}
Put $M=\sum_iX_i$ and take $0\le F\le1$ on configurations. For $q<p$,
use the symmetric test
$F_q(i_1,\ldots,i_q)=F(\sum_{v=1}^q\delta_{i_v})$ on ordered distinct
\emph{occurrences}; occurrences at the same location are still distinct.
After selecting $q$ occurrences, approximate the event that no others
remain by
\[
 \sum_{j=0}^{p-q-1}\frac{(-1)^j}{j!}(M-q)_j.
\]
For each integer $n\ge0$, the error in this approximation to
$\mathbf1_{\{n=0\}}$ is at most $(n)_{p-q}/(p-q)!$.
Indeed, when $n>0$ the partial binomial sum is
$(-1)^{p-q-1}\binom{n-1}{p-q-1}$; it vanishes when the cutoff exceeds $n$,
and otherwise its absolute value is at most $\binom{n}{p-q}$.

Divide by $q!$ to compensate for ordering the selected occurrences.
The error for this $q$ is at most
$\mu_p(I^p)/(q!(p-q)!)$. The omitted event $M\ge p$ costs at most
$\mu_p(I^p)/p!$. Summing over $q=0,\ldots,p-1$ bounds the remainder by
$2^p\mu_p(I^p)/p!$.
For the Poisson vector the analogous remainder is $2^p\lambda^p/p!$.
At total factorial order $j<p$, the sum of absolute coefficients is
\[
 \sum_{q=0}^j\frac1{q!(j-q)!}=\frac{2^j}{j!}.
\]
The order-zero terms agree. Subtraction of the two finite expansions and
supremization over indicator tests $F$ prove the lemma.
\end{proof}

\begin{proof}[Proof of Theorem~\ref{thm:bounded-mass}]
Choose $\alpha<\beta<2/3$. For every fixed $j$, all ordered lists of $j$
lengths from $I_d$ eventually have total at most $N^\beta$. Theorem~\ref{thm:moments}
gives a uniform error $\varepsilon_{d,j}\to0$ and hence
\[
 \|\mu_j-\pi_j\|_1\le\varepsilon_{d,j}\lambda_d^j,
 \qquad
 \mu_j(I_d^j)\le(1+\varepsilon_{d,j})\lambda_d^j.
\]
Substitution in \eqref{eq:10.2} yields, for fixed $p$,
\[
 d_{\rm TV}\le
 \sum_{j=1}^{p-1}\frac{(2\Lambda)^j}{j!}\varepsilon_{d,j}
 +\frac{(2\Lambda)^p}{p!}(2+\varepsilon_{d,p}).
 \tag{10.3}\label{eq:10.3}
\]
First let $d\to\infty$ with $p$ fixed, and then let $p\to\infty$.
No estimate uniform in growing moment order is required.
\end{proof}

\subsection{The scale-invariant cycle process}
The scale-invariant Poisson process has intensity $dx/x$ on
$(0,\infty)$; see Arratia~\cite{Arratia} for its role in permutation
cycles and other logarithmic structures. The next result identifies this
process as a moving mesoscopic limit for the present feedback model.

\begin{corollary}[Scale-invariant Poisson limit]\label{cor:scale-poisson}
Let $a_d/d\to\infty$, and assume $a_d\le N^\alpha$ for some fixed
$0<\alpha<2/3$. Then the random counting measures
\[
 \Xi_d=\sum_{r=1}^N C_r\,\delta_{r/a_d}
 \quad\Longrightarrow\quad
 \Xi\sim\operatorname{PPP}(dx/x)
 \tag{10.4}\label{eq:10.4}
\]
converge in distribution for the vague topology on locally finite counting
measures on $(0,\infty)$.
\end{corollary}
\begin{proof}
On a compact interval $[c,C]\subset(0,\infty)$, the relevant lengths
exceed $4d$ eventually and are at most $CN^\alpha$. Use a slightly larger
fixed exponent below $2/3$ in Theorem~\ref{thm:bounded-mass}. Their harmonic
mass is bounded, and the discrete measures
$\sum_r r^{-1}\delta_{r/a_d}$ converge vaguely to $dx/x$.
The Laplace functionals of the corresponding Poisson processes therefore
converge. Total variation transfers this convergence for compactly
supported tests. Bounded expected counts on each compact interval give
local tightness, proving the assertion.
\end{proof}

In particular, for $b=\lfloor N^\alpha\rfloor$ with $0<\alpha<2/3$,
\[
 \sum_{b<r\le2b}C_r\Longrightarrow\operatorname{Poisson}(\log2).
\]
A fixed finite collection of disjoint multiplicative windows has independent
limiting counts, with means equal to the logarithms of their endpoint ratios.
The windows may be at widely separated scales. These conclusions do not
apply to the macroscopic window $(N/2,N]$, which contains at most one cycle.

\subsection{Dickman mass and the largest cycle below a cutoff}
The relation between scale-invariant Poisson processes and Dickman laws is
classical \cite{Arratia,BM}. We give the short transfer proof to specify the
normalization and the Thorp application.

\begin{corollary}[Truncated cycle mass]\label{cor:dickman}
Fix $0<\alpha<2/3$, put $b=\lfloor N^\alpha\rfloor$, and define
\[
 T_d(b)=\frac1b\sum_{r\le b}rC_r.
\]
Then $T_d(b)\Longrightarrow D$, where the standard Dickman variable $D$ is
specified by
\[
 \mathbb E e^{-tD}
 =\exp\!\left\{\int_0^1\frac{e^{-tx}-1}{x}\,dx\right\},
 \qquad t\ge0.
 \tag{10.5}\label{eq:10.5}
\]
Equivalently, $D=\int_{(0,1]}x\,\Xi(dx)$ for the process in
Corollary~\ref{cor:scale-poisson}. The limiting mean and variance are
$\mathbb ED=1$ and $\operatorname{Var}(D)=1/2$.
\end{corollary}
\begin{proof}
For fixed $\varepsilon>0$, integration on $(\varepsilon,1]$ is continuous
at the limiting point process almost surely: it has no atoms at either
boundary. Thus the contribution from $\varepsilon b<r\le b$ converges by
Corollary~\ref{cor:scale-poisson}. The all-positive-length upper bound gives
\[
 \mathbb E\left[\frac1b\sum_{r\le\varepsilon b}rC_r\right]
 \le2\varepsilon
\]
for sufficiently large $d$. The limiting omitted Poisson mass has expectation
$\varepsilon$. Markov's inequality and $\varepsilon\downarrow0$ complete the
convergence argument. The Poisson Laplace functional gives
\eqref{eq:10.5}; its first two cumulants are
$\int_0^1x\,dx/x=1$ and $\int_0^1x^2\,dx/x=1/2$.
\end{proof}

\begin{corollary}[Largest cycle below a cutoff]\label{cor:maximum}
With $b$ as above, let
$R_d(b)=\max(\{r\le b:C_r>0\}\cup\{0\})$. Then
\[
 \frac{R_d(b)}b\Longrightarrow\operatorname{Uniform}(0,1).
 \tag{10.6}\label{eq:10.6}
\]
\end{corollary}
\begin{proof}
For $0<x<1$, the event $R_d(b)/b\le x$ is the absence of cycles in
$(xb,b]$. Its limiting probability, by Theorem~\ref{thm:bounded-mass}, is
$\exp(-\int_x^1dt/t)=x$.
\end{proof}

The harmonic mass of the entire range $[4d,N^\alpha]$ diverges. Consequently
Theorem~\ref{thm:bounded-mass} does not yield full-vector total-variation
approximation on that range. The next section obtains an explicit estimate
for an initial range below the birthday scale by a separate argument.

\section{Quantitative approximation below the birthday scale}\label{sec:birthday}
This section gives an explicit error for the complete initial cycle-count
vector. Its proof uses context-simple necklace indicators, not a growing-order
version of Theorem~\ref{thm:moments}.

\begin{theorem}[Quantitative initial-range approximation]\label{thm:birthday}
There is an absolute constant $C$ such that, for sufficiently large $d$ and
$2d\le b\le\sqrt N$,
\[
 d_{\rm TV}\!\left(
 \mathcal L((C_r)_{2d\le r\le b}),
 \bigotimes_{r=2d}^{b}\operatorname{Poisson}(1/r)
 \right)
 \le C d^4(\log_2d)^4\frac{b^2}{N}.
 \tag{11.1}\label{eq:11.1}
\]
In particular, the distance tends to zero when
$b=o(\sqrt N/[d^2(\log_2d)^2])$.
\end{theorem}

\subsection{Context-simple cycles and sequential forced closure}
A cycle is \emph{context-simple} if its queried length-$(d-1)$ contexts are
all distinct. Let $C_r^{\rm rep}$ count the remaining $r$-cycles, and let
$X_{r,s}$ count ordered pairs of distinct context-simple cycles of lengths
$r,s$ sharing at least one context. Equal lengths are allowed.

We first control one repeated-context cycle or a pair of context-simple
cycles. Throughout this section, use the following forced probability law:
choose independent uniform roots and construct each prescribed cycle by
$r-d$ ordinary steps and $d$ forced closing steps. For a pair, construct the
shorter cycle completely before constructing the longer one. Reject an
incompatible closure, a premature return, or an intersection of the cycles.
Let $V$ count closing queries already fixed by \emph{any} earlier query,
including earlier closing queries. If the completed family uses $v$ contexts,
the forced atom weight is $2^{-v-V}$. Thus, on a root-invariant class $E_m$
with $m$ repeated contexts and total length $L$, root averaging gives
\[
 \mathbb P^\circ(E_m)
 =\int_{E_m}\mathbb E_{\rm roots}(2^{-V})\,d\mu.
 \tag{11.2}\label{eq:11.2}
\]
Here $\mu$ is the corresponding scaled ordinary cycle measure.
For one cycle, $\mathbb E_{\rm roots}V\le2dm/r$. For a context-simple pair
$r\le s$, only the second cycle can have known closing queries, and
$\mathbb E_{\rm roots}V=dm/s\le2dm/(r+s)$. Therefore, writing $Q^{(m)}$
for the contribution of $E_m$,
\[
 Q^{(m)}\le2^{2dm/L}\mathbb P^\circ(E_m).
 \tag{11.3}\label{eq:11.3}
\]

Until the first repeated feedback query, this construction agrees with
independent fair cyclic seed words: the starting $d$ bits and the $r-d$
open-step output bits are independent, and closing appends the starting
bits. For $r\ge2d$, each pair of distinct $(d-1)$-contexts on a circle
matches with probability $2^{-(d-1)}=2/N$, by Appendix~\ref{app:contexts}.
For two different circles the same probability follows from independence.
On a successful intersecting context-simple pair, the first repetition
must be between the two circles. Hence
\[
 \begin{aligned}
 \sum_{m\ge1}\mathbb P^\circ(E_m)
 &\le r(r-1)/N &&\text{for one cycle},\\
 \sum_{m\ge1}\mathbb P^\circ(E_m)
 &\le 2rs/N &&\text{for a context-simple pair}.
 \end{aligned}
 \tag{11.4}\label{eq:11.4}
\]
No unexposed feedback value is resampled after its first query.

\subsection{A uniform sparse penalty}
For $d^4\le L\le2\sqrt N$, the one- or two-word type estimate is
\[
 \log_2 Q^{(m)}
 \le \frac{Lt}{h}
 +2M\log_2[e(1+L/M)]+2\log_2L-\kappa m,
 \quad M=2^h,
 \tag{11.5}\label{eq:11.5}
\]
where $t=\log_2L-h$ and $\kappa=(2d-3h)/h$.
Set $q=\log_2d$ and
\[
 h=\lfloor\log_2L-\log_2d-\log_2q\rfloor.
\]
For sufficiently large $d$, $2\le h\le d/2$, $\kappa\ge1$,
$L/(2dq)<M\le L/(dq)$, and
$2M\log_2[e(1+L/M)]\le4L/d$.
Put
\[
 F=Lt/h+4L/d+2\log_2L,
 \qquad m_0=\left\lceil\frac{F+L/d}{\kappa}\right\rceil.
\]
Then
\[
 \sum_{m\ge m_0}Q^{(m)}
 \le2^{1-L/d}\le2^{1-d^3}.
 \tag{11.6}\label{eq:11.6}
\]
For $m<m_0$, the penalty in \eqref{eq:11.3} is at most
$\Pi=2^{2dm_0/L}$. Direct expansion gives
\[
 \log_2\Pi
 \le\frac{2d}{2d-3h}
 \left[t+\frac{5h}{d}+\frac{2h\log_2L}{L}\right]+\frac{2d}{L}
 \le4t+15.
\]
Here the leading factor is at most $4$, $5h/d\le5/2$,
$2h\log_2L/L\le1$, and $2d/L\le1$.
Since $t<q+\log_2q+1$, one may take
\[
 \Pi\le 2^{20}d^4q^4.
 \tag{11.7}\label{eq:11.7}
\]
Combining \eqref{eq:11.3}--\eqref{eq:11.7} yields
\[
 \begin{aligned}
 \mathbb EC_r^{\rm rep}
 &\le\Pi\frac{r-1}{N}+\frac{2^{1-r/d}}r
 &&(r\ge d^4),\\
 \mathbb EX_{r,s}
 &\le\frac{2\Pi}{N}+\frac{2^{1-(r+s)/d}}{rs}
 &&(r+s\ge d^4).
 \end{aligned}
 \tag{11.8}\label{eq:11.8}
\]

For the remaining short ranges, the following elementary bounds suffice:
\[
 \mathbb EC_r^{\rm rep}\le\frac{r(r-1)}N,
 \qquad
 \mathbb EX_{r,s}\le\frac{4\sqrt{rs}}N.
 \tag{11.9}\label{eq:11.9}
\]
For one cycle, $H_d=\log_2r$ and $H_d-H_{d-1}=2m/r$ imply
$2dm/r\le\log_2r$. Jensen in \eqref{eq:11.2} therefore loses at most a
factor $r$; \eqref{eq:11.4} proves the first bound.
For a context-simple pair, choose one of the two circles with probability
$1/2$, then a position uniformly on that circle. This stationary mixture
has $H_d=1+\tfrac12\log_2(rs)$. Each shared context merges atoms of masses
$1/(2r)$ and $1/(2s)$ and contributes at least $1/s$ to the final entropy
increment. Thus $dm/s\le1+\tfrac12\log_2(rs)$, and the inverse root-averaged
likelihood is at most $2\sqrt{rs}$. The second bound follows from
\eqref{eq:11.4} after dividing by $rs$.

Summing \eqref{eq:11.8} and \eqref{eq:11.9} gives
\[
 \sum_{r=2d}^b\mathbb EC_r^{\rm rep}
 +\sum_{r,s=2d}^b\mathbb EX_{r,s}
 \le C d^4q^4\frac{b^2}{N}.
 \tag{11.10}\label{eq:11.10}
\]
Indeed, on $r<d^4$ use $r^2\le d^4r$, and on $r+s<d^4$ use
$\sqrt{rs}\le d^4/2$. The dense remainders are at most a polynomial in $b$
times $2^{-d^3}$ and are negligible uniformly for $b\le\sqrt N$.

\subsection{The point-process comparison}
Let $\mathcal H_r$ be the set of possible context-simple primitive necklaces
of length $r$. The event that $\gamma\in\mathcal H_r$ occurs as a cycle has
indicator $I_\gamma$ and probability $p_\gamma=2^{-r}$.
Its dependency neighborhood consists of all potential necklaces sharing
one of its contexts, including itself. The indicator is independent of the
joint collection outside that neighborhood, since their determining
feedback bits are disjoint.

The point-process theorem of Arratia, Goldstein, and Gordon
\cite[Theorem~2]{AGG}, in our total-variation convention, bounds the distance
from the independent Poisson indicators by $2(b_1+b_2)$, where
\[
 b_1=\sum_\gamma\sum_{\gamma'\sim\gamma}p_\gamma p_{\gamma'},
 \qquad
 b_2=\sum_\gamma\sum_{\substack{\gamma'\sim\gamma\\\gamma'\ne\gamma}}
 \mathbb E(I_\gamma I_{\gamma'}).
\]
The external-dependence term is zero.
For each pair of lengths $r,s$, rooting both necklaces gives
\[
 \sum_{\substack{\gamma\in\mathcal H_r,\ \gamma'\in\mathcal H_s\\
                  \gamma\sim\gamma'}}p_\gamma p_{\gamma'}
 \le\frac1{rs}\mathbb P\{
 \text{two independent seed words share a context}\}
 \le\frac2N.
\]
This includes the diagonal $\gamma=\gamma'$.
Consequently $b_1\le2b^2/N$. Simultaneously occurring distinct necklaces
are distinct actual cycles, so \eqref{eq:11.10} gives
$b_2\le\sum_{r,s}\mathbb EX_{r,s}=O(d^4q^4b^2/N)$.
Aggregation by length cannot increase total variation.

Let $\lambda_r=|\mathcal H_r|2^{-r}$ be the mean context-simple cycle count.
A length-$r$ cyclic seed has all contexts distinct if and only if it is a
rooted context-simple necklace. Therefore
\[
 r\lambda_r=1-\delta_d(r),
 \qquad
 0\le\frac1r-\lambda_r\le\frac{r-1}{N},
\]
where $\delta_d(r)\le r(r-1)/N$ is its context-collision probability.
Changing the independent Poisson means to $1/r$ costs $O(b^2/N)$.
Restoring non-context-simple cycles costs at most
$\sum_r\mathbb EC_r^{\rm rep}$, bounded in \eqref{eq:11.10}.
This proves Theorem~\ref{thm:birthday}.

Theorems~\ref{thm:bounded-mass} and \ref{thm:birthday} are complementary.
The former applies beyond the birthday scale but only to sets of bounded
harmonic mass, with no asserted rate; the latter controls the complete
initial vector and gives an explicit rate. Neither asserts a Poisson
approximation for that complete vector above the birthday scale.

\section{Scope and further directions}\label{sec:scope}
The parameter restrictions are fixed before $d\to\infty$.
Theorem~\ref{thm:moments} is uniform in length ratios but not in a growing
number of cycles. The functional theorem concerns cycles of length at most
$N^\alpha$, for fixed $\alpha<2/3$, rather than the full order of $T_f$.
The Poisson theorem on bounded-harmonic-mass sets does allow a growing number
of coordinates, but it does not cover the whole initial vector when its
harmonic mass diverges.

The full one-shuffle Erd\H{o}s--Tur\'{a}n statement suggested by these
results is
\[
 \frac{\log\operatorname{ord}(T_f)-\tfrac12(\log N)^2}
 {\sqrt{(\log N)^3/3}}
 \Longrightarrow\mathcal N(0,1).
\]
It is not proved here. The present dense-tail estimate requires
$\kappa=(2d-3h)/h>0$. Extending the moment comparison to the remaining
length scales therefore requires an additional estimate, not just taking
$\beta$ to the endpoint in a fixed-parameter theorem.

The lower and upper moment arguments are analytic and combinatorial.
Finite integer-arithmetic checks of feedback switches, root weights,
entropy inequalities, and bridge identities accompany the source. They
are reproducibility aids, not inputs to the asymptotic proof. The separate
repository documentation records the checks, review history, and their
limitations.

\appendix
\section{Equality graphs and exact short cycles}\label{app:contexts}
\subsection{Wrapped-context comparisons}

This appendix proves the small-length collision estimate used in \eqref{eq:6.1}.
Let \(B\) require primitive component words and different necklaces for
any two components of the same length.

For two positions on the same circle of length \(n\), with nonzero
displacement \(a\), the \(h\) consecutive comparisons are edges
\(t\leftrightarrow t+a\) modulo \(n\). Write \(g=\gcd(n,a)\). For
\(h<n\), the omitted \(n-h\) starting positions are consecutive and meet
\(\min(g,n-h)\) of the \(g\) cycles of the full shift graph. The
retained equality system therefore has rank \[
 h-\max(0,g-n+h)=\min(h,n-g)
\] over \(\mathbb F_2\). The same formula holds for \(h\ge n\), when the
rank has already saturated. If \(h\ge n-g\), equality forces period
dividing \(g<n\), which is impossible on \(B\). Otherwise the rank is
\(h\), so equality has probability \(2^{-h}\). This includes a two-cycle
in the shift graph: its two oriented edges impose one linear relation.

For positions on different circles of lengths \(r,s\), offset their
roots so that the comparisons are
\((t\bmod r)\leftrightarrow(t\bmod s)\), \(0\le t<h\), in a bipartite
graph. Let \(g=\gcd(r,s)\) and assume \(r\le s\). Its first \(s\) edges
attach each vertex of the second circle once. Contract those edges. The
next \(r-g\) edges become \(j\leftrightarrow j+s\pmod r\),
\(0\le j<r-g\). The full shift-by-\(s\) graph has \(g\) cycles, and the
omitted \(g\) consecutive edges remove exactly one from each. Thus the
first \(r+s-g\) comparison edges form a forest and have that rank. At
this threshold both words are forced to have period dividing \(g\). This
is excluded by \(B\): either a component is nonprimitive, or the two
have equal length and represent the same necklace. For smaller \(h\) the
comparisons have rank \(h\) and probability \(2^{-h}\).

Taking \(h=d-1\) proves, for every pair of distinct positions, \[
 \mathbb P\{B\text{ and their contexts agree}\}\le2^{-(d-1)}.
\] On \(B\) without such a collision, the weighted-word integrand equals
one. On valid words its maximum is \(2^m\le L^{L/(2d)}\le L^{2k}\) when
every length is below \(4d\). The union bound over \(\binom L2\) pairs
gives \eqref{eq:6.1}.

\subsection{Exact short-cycle laws}
Let
\[
 \eta_r=\frac1r\sum_{a\mid r}\mu(a)2^{r/a}
\]
be the number of primitive binary necklaces of length $r$, where $\mu$ is
the M\"obius function. The independent-binomial phenomenon was proved by
Stark~\cite{Stark}. The wrapped-context comparison above yields the following
explicit range directly.

\begin{proposition}[Independent short-cycle counts]\label{prop:short-cycles}
For $1\le r\le d-1$,
\[
 C_r\sim\operatorname{Binomial}(\eta_r,2^{-r}).
\]
For $d\ge2$, the variables
$C_1,\ldots,C_{\lfloor(d+1)/2\rfloor}$ are mutually independent.
\end{proposition}
\begin{proof}
A primitive necklace of length $r\le d-1$ has $r$ distinct contexts and
occurs with probability $2^{-r}$. Distinct necklaces of the same length
have disjoint context sets: agreement on $d-1\ge r$ consecutive positions
would identify their periodic extensions. This gives the binomial marginal.

Put $b=\lfloor(d+1)/2\rfloor$. For different lengths $r,s\le b$,
\[
 r+s-\gcd(r,s)\le2b-2\le d-1.
\]
The two-circle comparison above shows that a common context would identify
the periodic extensions, contradicting their distinct minimal periods.
Thus every primitive-necklace indicator involved in these lengths depends
on a disjoint set of feedback bits. This proves joint independence.
\end{proof}

Stark's stated joint range is $r\le(d-1)/3$; the proposition records the
larger elementary range furnished by the present equality-graph argument.
In particular $C_1\sim\operatorname{Binomial}(2,1/2)$, so a Poisson
approximation to all fixed small lengths cannot hold. This explains the
lower cutoff in the harmonic moment asymptotics.

\begin{thebibliography}{9}

\bibitem{Arratia}
R.~Arratia,
\emph{On the central role of scale invariant Poisson processes on $(0,\infty)$},
in \emph{Microsurveys in Discrete Probability} (Princeton, NJ, 1997),
DIMACS Series in Discrete Mathematics and Theoretical Computer Science,
vol.~41, American Mathematical Society, Providence, RI, 1998, pp.~21--41.
Updated electronic version: \href{https://arxiv.org/abs/1611.05572}{arXiv:1611.05572}.

\bibitem{ACH}
R.~A.~Arratia, E.~R.~Canfield, and A.~W.~Hales,
\emph{Random feedback shift registers and the limit distribution for largest cycle lengths},
Combinatorics, Probability and Computing \textbf{32} (2023), no.~4, 559--593.
\href{https://doi.org/10.1017/S0963548323000020}{doi:10.1017/S0963548323000020}.

\bibitem{AGG}
R.~Arratia, L.~Goldstein, and L.~Gordon,
\emph{Poisson approximation and the Chen--Stein method},
Statistical Science \textbf{5} (1990), no.~4, 403--424.
\href{https://doi.org/10.1214/ss/1177012015}{doi:10.1214/ss/1177012015}.

\bibitem{BM}
C.~Bhattacharjee and I.~Molchanov,
\emph{Convergence to scale-invariant Poisson processes and applications in Dickman approximation},
Electronic Journal of Probability \textbf{25} (2020), paper no.~79, 1--20.
\href{https://doi.org/10.1214/20-EJP482}{doi:10.1214/20-EJP482}.

\bibitem{ET}
P.~Erd\H{o}s and P.~Tur\'{a}n,
\emph{On some problems of a statistical group-theory. I},
Zeitschrift f\"ur Wahrscheinlichkeitstheorie und Verwandte Gebiete
\textbf{4} (1965), 175--186.
\href{https://doi.org/10.1007/BF00536750}{doi:10.1007/BF00536750}.

\bibitem{Long}
C.~D.~Long,
\emph{Erd\H{o}s--Tur\'{a}n universality after one shuffle},
research manuscript, 2026. Source and PDF in the
\href{https://github.com/long-mathematics/nonuniform-random-permutations}{Nonuniform Random Permutations repository}, Paper~V.

\bibitem{Stark}
D.~Stark,
\emph{The small cycle counts of random feedback shift registers},
Australasian Journal of Combinatorics \textbf{86} (2023), no.~3, 414--422.
\url{https://ajc.maths.uq.edu.au/pdf/86/ajc_v86_p414.pdf}.

\bibitem{Thorp}
E.~O.~Thorp,
\emph{Nonrandom shuffling with applications to the game of Faro},
Journal of the American Statistical Association \textbf{68} (1973),
no.~344, 842--847.
\href{https://doi.org/10.1080/01621459.1973.10481434}{doi:10.1080/01621459.1973.10481434}.

\end{thebibliography}

\end{document}
